\PassOptionsToPackage{draft}{hyperref}
\documentclass[12pt,a4paper]{article}

\usepackage{viu-mrob-thesis}
\usepackage{csquotes}
\usepackage{multicol}
\usepackage{needspace}
\usepackage[
  backend=biber,
  style=apa,
  natbib=true,
  hyperref=false,
  sorting=nyt,
  sortcites=true
]{biblatex}
\addbibresource{references.bib}
\renewcommand*{\bibfont}{\raggedright}

\newcommand{\SPTwoReadyDynamicRuns}{2100}
\newcommand{\SPTwoReadyWorlds}{210}
\newcommand{\SPTwoReadyNOneRows}{36600}
\newcommand{\SPTwoReadyNOneFalse}{4536}
\newcommand{\SPTwoReadyNOneFalseRate}{0{,}124}
\newcommand{\SPTwoReadyNTwoRows}{720}
\newcommand{\SPTwoReadyNTwoFalse}{52}
\newcommand{\SPTwoReadyNTwoRejectRate}{7{,}222}
\newcommand{\SPTwoReadyNTwoFeasible}{668}
\newcommand{\SPTwoReadyNThreeRows}{240}
\newcommand{\SPTwoReadyNThreeLeaderBytes}{50478}
\newcommand{\SPTwoReadyNThreeVirtualBytes}{192876}
\newcommand{\SPTwoReadyNThreeLeaderPositionRMSE}{0{,}182}
\newcommand{\SPTwoReadyNThreeVirtualPositionRMSE}{0{,}13}
\newcommand{\SPTwoReadyNThreeLeaderYawRMSE}{0{,}04}
\newcommand{\SPTwoReadyNThreeVirtualYawRMSE}{0{,}026}
\newcommand{\SPTwoReadyNThreeLossyPositionDifference}{0{,}132}
\newcommand{\SPTwoReadyNThreeLossyPositionCILow}{0{,}131}
\newcommand{\SPTwoReadyNThreeLossyPositionCIHigh}{0{,}134}
\newcommand{\SPTwoReadyNThreeLossyLeaderAge}{0{,}105}
\newcommand{\SPTwoReadyNThreeLossyVirtualAge}{1{,}971}
\newcommand{\SPTwoReadyNThreeLossyLeaderDelivery}{0{,}936}
\newcommand{\SPTwoReadyNThreeLossyVirtualDelivery}{0{,}937}
\newcommand{\SPTwoReadyNThreeLeaderDeliveredBytes}{44596}
\newcommand{\SPTwoReadyNThreeVirtualDeliveredBytes}{185017}
\newcommand{\SPTwoReadyGovernedSuccess}{0{,}794}
\newcommand{\SPTwoReadyRawSuccess}{0{,}74}
\newcommand{\SPTwoReadyGovernedCollision}{0{,}1}
\newcommand{\SPTwoReadyRawCollision}{0{,}036}
\newcommand{\SPTwoReadySSixSuccess}{0}
\newcommand{\SPTwoReadySThreeSuccess}{0{,}628}
\newcommand{\SPTwoReadySFiveSuccess}{1}
\newcommand{\SPTwoReadySSixCollision}{0{,}7}
\newcommand{\SPTwoReadyPDGovernedSuccess}{0{,}8}
\newcommand{\SPTwoReadyPDRawSuccess}{0{,}724}
\newcommand{\SPTwoReadyLQIGovernedSuccess}{0{,}771}
\newcommand{\SPTwoReadyLQIRawSuccess}{0{,}657}
\newcommand{\SPTwoReadyPHGovernedSuccess}{0{,}8}
\newcommand{\SPTwoReadyPHRawSuccess}{0{,}724}
\newcommand{\SPTwoReadyMPCGovernedSuccess}{0{,}857}
\newcommand{\SPTwoReadyMPCRawSuccess}{0{,}857}
\newcommand{\SPTwoReadyPDDifference}{0{,}076}
\newcommand{\SPTwoReadyPDCILow}{0{,}048}
\newcommand{\SPTwoReadyPDCIHigh}{0{,}105}
\newcommand{\SPTwoReadyPDHolmP}{\ensuremath{<0{,}001}}
\newcommand{\SPTwoReadyPDBlockHolmP}{\ensuremath{<0{,}001}}
\newcommand{\SPTwoReadyLQIDifference}{0{,}114}
\newcommand{\SPTwoReadyLQICILow}{0{,}081}
\newcommand{\SPTwoReadyLQICIHigh}{0{,}148}
\newcommand{\SPTwoReadyLQIHolmP}{\ensuremath{<0{,}001}}
\newcommand{\SPTwoReadyLQIBlockHolmP}{\ensuremath{<0{,}001}}
\newcommand{\SPTwoReadyPHDifference}{0{,}076}
\newcommand{\SPTwoReadyPHCILow}{0{,}052}
\newcommand{\SPTwoReadyPHCIHigh}{0{,}1}
\newcommand{\SPTwoReadyPHHolmP}{\ensuremath{<0{,}001}}
\newcommand{\SPTwoReadyPHBlockHolmP}{\ensuremath{<0{,}001}}
\newcommand{\SPTwoReadyLQICollisionDifference}{0{,}114}
\newcommand{\SPTwoReadyMPCDifference}{0}
\newcommand{\SPTwoReadyMPCCILow}{0}
\newcommand{\SPTwoReadyMPCCIHigh}{0}
\newcommand{\SPTwoReadyMPCHolmP}{1{,}000}
\newcommand{\SPTwoReadyMPCBlockHolmP}{1{,}000}
\newcommand{\SPTwoReadyPDCollisionDifference}{0{,}124}
\newcommand{\SPTwoReadyPDCollisionCILow}{0{,}105}
\newcommand{\SPTwoReadyPDCollisionCIHigh}{0{,}138}
\newcommand{\SPTwoReadyLQICollisionCILow}{0{,}095}
\newcommand{\SPTwoReadyLQICollisionCIHigh}{0{,}133}
\newcommand{\SPTwoReadyPHCollisionDifference}{0{,}114}
\newcommand{\SPTwoReadyPHCollisionCILow}{0{,}09}
\newcommand{\SPTwoReadyPHCollisionCIHigh}{0{,}133}
\newcommand{\SPTwoReadyMPCCollisionDifference}{-0{,}143}
\newcommand{\SPTwoReadyMPCCollisionCILow}{-0{,}143}
\newcommand{\SPTwoReadyMPCCollisionCIHigh}{-0{,}143}
\newcommand{\SPTwoReadyPairedCollisionDifference}{0{,}052}
\newcommand{\SPTwoReadyPairedCollisionCILow}{0{,}039}
\newcommand{\SPTwoReadyPairedCollisionCIHigh}{0{,}064}
\newcommand{\SPTwoReadyPairedSuccessDifference}{0{,}067}
\newcommand{\SPTwoReadyPairedSuccessCILow}{0{,}05}
\newcommand{\SPTwoReadyPairedSuccessCIHigh}{0{,}082}
\newcommand{\SPTwoReadyPairedAdmissibleSuccessDifference}{0{,}41}
\newcommand{\SPTwoReadyPairedAdmissibleSuccessCILow}{0{,}38}
\newcommand{\SPTwoReadyPairedAdmissibleSuccessCIHigh}{0{,}437}
\newcommand{\SPTwoReadyPDRawPhysicalSuccess}{0}
\newcommand{\SPTwoReadyPDGovernedPhysicalSuccess}{0{,}429}
\newcommand{\SPTwoReadyPDRawViolation}{1}
\newcommand{\SPTwoReadyPDGovernedViolation}{0{,}571}
\newcommand{\SPTwoReadyPDRawCommandInfeasible}{1}
\newcommand{\SPTwoReadyPDGovernedCommandInfeasible}{0{,}286}
\newcommand{\SPTwoReadyLQIRawPhysicalSuccess}{0}
\newcommand{\SPTwoReadyLQIGovernedPhysicalSuccess}{0{,}352}
\newcommand{\SPTwoReadyLQIRawViolation}{1}
\newcommand{\SPTwoReadyLQIGovernedViolation}{0{,}638}
\newcommand{\SPTwoReadyLQIRawCommandInfeasible}{1}
\newcommand{\SPTwoReadyLQIGovernedCommandInfeasible}{0{,}286}
\newcommand{\SPTwoReadyPHRawPhysicalSuccess}{0}
\newcommand{\SPTwoReadyPHGovernedPhysicalSuccess}{0{,}41}
\newcommand{\SPTwoReadyPHRawViolation}{1}
\newcommand{\SPTwoReadyPHGovernedViolation}{0{,}59}
\newcommand{\SPTwoReadyPHRawCommandInfeasible}{1}
\newcommand{\SPTwoReadyPHGovernedCommandInfeasible}{0{,}286}
\newcommand{\SPTwoReadyMPCRawPhysicalSuccess}{0}
\newcommand{\SPTwoReadyMPCGovernedPhysicalSuccess}{0{,}448}
\newcommand{\SPTwoReadyMPCRawViolation}{1}
\newcommand{\SPTwoReadyMPCGovernedViolation}{0{,}552}
\newcommand{\SPTwoReadyMPCRawCommandInfeasible}{1}
\newcommand{\SPTwoReadyMPCGovernedCommandInfeasible}{0{,}286}
\newcommand{\SPTwoReadyPairedRawRealizedFraction}{0{,}507}
\newcommand{\SPTwoReadyPairedGovernedRealizedFraction}{0{,}135}
\newcommand{\SPTwoReadyPairedRealizedFractionDifferenceCILow}{-0{,}378}
\newcommand{\SPTwoReadyPairedRealizedFractionDifferenceCIHigh}{-0{,}366}
\newcommand{\SPTwoReadyPairedRawMaxViolationDuration}{3{,}361}
\newcommand{\SPTwoReadyPairedGovernedMaxViolationDuration}{1{,}038}
\newcommand{\SPTwoReadyPairedDurationDifferenceCILow}{-2{,}444}
\newcommand{\SPTwoReadyPairedDurationDifferenceCIHigh}{-2{,}196}
\newcommand{\SPTwoReadyPairedRawMinKinematicMargin}{-3{,}811}
\newcommand{\SPTwoReadyPairedGovernedMinKinematicMargin}{-0{,}188}
\newcommand{\SPTwoReadyPairedMinKinematicDifferenceCILow}{2{,}868}
\newcommand{\SPTwoReadyPairedMinKinematicDifferenceCIHigh}{4{,}489}
\newcommand{\SPTwoReadyPairedRawMinMechanicalMargin}{0{,}547}
\newcommand{\SPTwoReadyPairedGovernedMinMechanicalMargin}{0{,}589}
\newcommand{\SPTwoReadyPDStepTenSuccess}{0{,}8}
\newcommand{\SPTwoReadyPDStepFiveSuccess}{0{,}8}
\newcommand{\SPTwoReadyLQIStepTenSuccess}{0{,}781}
\newcommand{\SPTwoReadyLQIStepFiveSuccess}{0{,}771}
\newcommand{\SPTwoReadyPHStepTenSuccess}{0{,}767}
\newcommand{\SPTwoReadyPHStepFiveSuccess}{0{,}786}
\newcommand{\SPTwoReadyMPCStepTenSuccess}{0{,}857}
\newcommand{\SPTwoReadyMPCStepFiveSuccess}{0{,}857}
\newcommand{\SPTwoReadyPDStepFivePhysicalSuccess}{0{,}167}
\newcommand{\SPTwoReadyPDStepFiveViolation}{0{,}833}
\newcommand{\SPTwoReadyLQIStepFivePhysicalSuccess}{0{,}195}
\newcommand{\SPTwoReadyLQIStepFiveViolation}{0{,}795}
\newcommand{\SPTwoReadyPHStepFivePhysicalSuccess}{0{,}143}
\newcommand{\SPTwoReadyPHStepFiveViolation}{0{,}857}
\newcommand{\SPTwoReadyMPCStepFivePhysicalSuccess}{0{,}224}
\newcommand{\SPTwoReadyMPCStepFiveViolation}{0{,}776}

\newcommand{\SPTwoNOneRows}{284}
\newcommand{\SPTwoNOneFalse}{14}
\newcommand{\SPTwoNOneSmallOmega}{0{,}56}
\newcommand{\SPTwoNOneLargeOmega}{0{,}16}
\newcommand{\SPTwoNTwoCounterexamples}{1}
\newcommand{\SPTwoNTwoRows}{288}
\newcommand{\SPTwoNTwoFeasible}{276}
\newcommand{\SPTwoNTwoInfeasible}{12}
\newcommand{\SPTwoNTwoResidual}{1{,}07\times 10^{-14}}
\newcommand{\SPTwoNThreeLeaderMessages}{60}
\newcommand{\SPTwoNThreeVirtualMessages}{160}
\newcommand{\SPTwoNThreeLeaderBytes}{8190}
\newcommand{\SPTwoNThreeVirtualBytes}{35600}
\newcommand{\SPTwoCagingClosed}{sin escape a la resolución fijada}
\newcommand{\SPTwoCagingGap}{escape encontrado}

\newcommand{\CargoEtwoEWorlds}{360}
\newcommand{\CargoEtwoERuns}{2160}
\newcommand{\CargoEtwoESuccess}{0{,}997}
\newcommand{\CargoEtwoEFailureSuccess}{0{,}989}
\newcommand{\CargoEtwoEDegradedSuccess}{1}
\newcommand{\CargoEtwoEMessages}{328{,}1}


\renewcommand{\viuownsource}{%
  \par\nobreak\vspace{0.8mm}\noindent
  \makebox[\linewidth][c]{\footnotesize\itshape Elaboración propia.}%
  \par
}

\setviuheadertitle{SP2 · Ejecución y transporte cooperativo}
\hypersetup{
  pdftitle={SP2: ejecución y transporte cooperativo con Pioneer 3-DX},
  pdfauthor={Jorge Luis Mayorga Taborda},
  pdfsubject={Rigidez cinemática, fuerzas, información local y validez física}
}

\newcommand{\leveltag}[1]{%
  \Needspace{5.5\baselineskip}%
  \noindent\colorbox{VIUDark}{%
    \parbox{\dimexpr\textwidth-2\fboxsep\relax}{%
      \color{white}\sffamily\bfseries\small #1%
    }%
  }\par\smallskip
}

\newcommand{\resultfigure}[4][0.40\textheight]{%
  \begin{center}
    \includegraphics[width=\textwidth,height=#1,keepaspectratio]{#2}
    \captionof{figure}{#3}
    \label{#4}
    \viuownsource
  \end{center}
}

% Silueta cenital reutilizable del Pioneer 3-DX. Los cinco argumentos son
% coordenadas, orientación, identificador y escala. El triángulo marca el frente.
\newcommand{\pioneericon}[5]{%
  \begin{scope}[shift={(#1,#2)},rotate=#3,scale=#5]
    \fill[VIUDark!88] (-.24,-.38) rectangle (.22,-.30);
    \fill[VIUDark!88] (-.24,.30) rectangle (.22,.38);
    \draw[fill=SPBlue!10,draw=SPBlue,line width=.75pt,rounded corners=.08cm]
      (-.42,-.28) rectangle (.34,.28);
    \draw[fill=SPBlue,draw=SPBlue] (.32,-.17)--(.54,0)--(.32,.17)--cycle;
    \draw[SPBlue!70!black,line width=.45pt]
      (.23,-.13) arc[start angle=-58,end angle=58,radius=.23];
    \node[font=\sffamily\bfseries\tiny,text=VIUDark] at (-.07,0) {#4};
  \end{scope}%
}

\newtheorem{sp2prop}{Proposición}
\newtheorem{sp2lema}{Lema}

\ExplSyntaxOn
\NewDocumentCommand{\spdecimal}{m}
  {
    \tl_set:Nx \l_tmpa_tl {#1}
    \tl_replace_all:Nnn \l_tmpa_tl {.} {{,}}
    \tl_use:N \l_tmpa_tl
  }
\ExplSyntaxOff

\begin{document}

\viuromanfoliofalse
\pagenumbering{arabic}
\setcounter{page}{1}
% Mantiene el interlineado VIU de 1,5 y recupera espacio mediante separaciones
% editoriales, no reduciendo la tipografía ni la escala de la Figura 1.
\setlength{\parskip}{1.5pt plus .5pt}
\setlength{\abovedisplayskip}{5pt plus 1pt minus 1pt}
\setlength{\belowdisplayskip}{5pt plus 1pt minus 1pt}
\setlength{\abovedisplayshortskip}{3pt plus 1pt}
\setlength{\belowdisplayshortskip}{3pt plus 1pt}
\setlength{\textfloatsep}{8pt plus 2pt minus 2pt}
\setlength{\intextsep}{8pt plus 2pt minus 2pt}

% --------------------------------------------------------------------------
% Página 1 · Introducción y frontera física
% --------------------------------------------------------------------------
\leveltag{CARGA RÍGIDA Y CONFINAMIENTO GEOMÉTRICO}
\section*{SP2: ejecución y transporte cooperativo}

El objetivo de SP2 no termina en decidir si una coalición es factible. SP1
entrega los miembros; SP2 estudia cómo llevar la carga hasta una pose sin
rebasar restricciones de movilidad, contacto e información. La campaña evalúa
una versión acotada, con geometría de contacto fija y planta plana, que escala
la aceleración nominal $a_L^0$ mediante $\alpha\in[0,1]$ antes de aplicarla:

\[
 a_L=\alpha a_L^0,\qquad
 \xi_L^+\in\Xi_{\mathcal C}^{\mathrm{align}},\qquad
 W_d(a_L)\in\widehat{\mathcal W}_{\mathcal C},\qquad
 g_{\mathrm{info}}=1 .
\]

En las cinco leyes locales, $g_{\mathrm{info}}$ solo controla la edad de la
estimación del AMR de referencia. No informa sobre la frescura conjunta de la
coalición. El MPC central usa estado global y fija esta guarda a uno.

Una orden eficaz en un modelo sin contacto puede quedar bloqueada al añadir
estas comprobaciones, lo que modifica la dinámica de lazo cerrado. El efecto se
mide con contrastes pareados dentro de cada familia. Promediar reguladores
distintos mezclaría modelos y sintonías.

En \textbf{Cargo}, modo primario, cada AMR sostiene la carga mediante un
acoplamiento que transmite fuerza planar y permite el giro propio del chasis.
En \textbf{caging}, extensión no prehensil, contactos unilaterales mantienen el
confinamiento; sus normales $f_{n,i}^{E}\geq0$ producen el wrench ilustrativo
$W_\Sigma^{E}=(F_\Sigma^{E},\tau_\Sigma^{E})$. La
Figura~\ref{fig:sp2-frontera} distingue las condiciones de cierre y fuerza
\parencite{bicchi1995closure,pereira2004caging}.

\clearpage
\begin{center}
\resizebox{.98\linewidth}{!}{%
\begin{tikzpicture}[x=1cm,y=1cm,>=Latex,
  font=\fontsize{9.6}{11}\selectfont\sffamily]
  \tikzset{
    modepanel/.style={rounded corners=2mm,fill=SPFloor,draw=SPGrid,line width=.7pt},
    scene/.style={rounded corners=1.2mm,fill=white,draw=SPGrid,line width=.6pt},
    stage/.style={rounded corners=1.2mm,fill=white,draw=SPGrid,line width=.65pt,
      minimum width=2.18cm,text width=1.96cm,minimum height=.90cm,
      align=center,font=\fontsize{7.8}{8.8}\selectfont,inner sep=2.2pt},
    flow/.style={->,line width=1.1pt,draw=VIUOrange,
      shorten >=2.2pt,shorten <=2.2pt},
    unilateral/.style={->,line width=1.35pt,draw=VIUOrange},
    contactforce/.style={->,draw=SPRed,shorten >=.5pt,line cap=round},
    resultant/.style={->,line width=1.45pt,draw=SPBlue,line cap=round},
    motion/.style={->,line width=1.35pt,draw=SPBlue},
    factcard/.style={rounded corners=1mm,draw=SPGrid,line width=.6pt,align=center}
  }

  % Dos secuencias temporales, cada una con estado inicial y final.
  \draw[modepanel] (0,2.26) rectangle (7.15,9.92);
  \draw[modepanel] (7.85,2.26) rectangle (15.00,9.92);
  \fill[SPBlue!9,rounded corners=2mm] (0,8.82) rectangle (7.15,9.92);
  \fill[VIUOrange!10,rounded corners=2mm] (7.85,8.82) rectangle (15.00,9.92);
  \node[anchor=west,font=\bfseries,text=VIUDark] at (.30,9.53)
    {CARGO · TRANSPORTE RÍGIDO};
  \node[anchor=west,font=\fontsize{8}{9.2}\selectfont,text=SPBlue!80!black] at (.30,9.08)
    {pivotes solidarios; chasis con giro libre};
  \node[anchor=west,font=\bfseries,text=VIUDark] at (8.15,9.53)
    {CAGING · EMPUJE CONFINADO};
  \node[anchor=west,font=\fontsize{8}{9.2}\selectfont,text=VIUOrange!85!black] at (8.15,9.08)
    {cerco reconfigurable sin vía de escape};

  % Cuatro ventanas con la misma retícula: dos estados por modalidad.
  \foreach \xa/\xb in {.18/3.34,3.81/6.97,8.03/11.19,11.66/14.82}{
    \draw[scene] (\xa,3.62) rectangle (\xb,8.62);
  }
  \foreach \x in {.92,1.76,2.60,4.55,5.39,6.23,8.77,9.61,10.45,
                   12.40,13.24,14.08}{
    \draw[SPGrid!45,line width=.28pt] (\x,3.82)--(\x,8.16);
  }
  \foreach \y in {4.54,5.58,6.62,7.66}{
    \draw[SPGrid!45,line width=.28pt] (.34,\y)--(3.18,\y);
    \draw[SPGrid!45,line width=.28pt] (3.97,\y)--(6.81,\y);
    \draw[SPGrid!45,line width=.28pt] (8.19,\y)--(11.03,\y);
    \draw[SPGrid!45,line width=.28pt] (11.82,\y)--(14.66,\y);
  }
  \node[anchor=west,font=\fontsize{8}{9.2}\selectfont\bfseries,text=VIUDark] at (.38,8.35)
    {(a) ANTES · $t_0$};
  \node[anchor=west,font=\fontsize{8}{9.2}\selectfont\bfseries,text=SPBlue!80!black] at (4.01,8.35)
    {(b) DESPUÉS · $t_f$};
  \node[anchor=west,font=\fontsize{8}{9.2}\selectfont\bfseries,text=VIUDark] at (8.23,8.35)
    {(c) ANTES · $t_0$};
  \node[anchor=west,font=\fontsize{8}{9.2}\selectfont\bfseries,text=VIUOrange!85!black] at (11.86,8.35)
    {(d) DESPUÉS · $t_f$};
  \node[font=\bfseries,text=SPBlue] at (3.58,6.02) {$\Longrightarrow$};
  \node[font=\bfseries,text=VIUOrange] at (11.42,6.02) {$\Longrightarrow$};

  % El contenido de las cuatro escenas queda centrado en la nueva altura útil.
  \begin{scope}[yshift=17.5mm]
  % (a) Cargo inicial: cuatro plataformas bajo las esquinas y rumbo común.
  \pioneericon{.88}{3.88}{14}{}{.43}
  \pioneericon{2.24}{3.88}{14}{}{.43}
  \pioneericon{.88}{4.66}{14}{}{.43}
  \pioneericon{2.24}{4.66}{14}{}{.43}
  \draw[rounded corners=.7mm,fill=VIUOrange!18,fill opacity=.72,
    draw=VIUOrange,draw opacity=1,line width=1pt]
    (.86,3.87) rectangle (2.25,4.67);
  \node[font=\bfseries\fontsize{8}{9.2}\selectfont] at (1.56,4.28) {carga $\mathcal L$};
  \foreach \p in {(.88,3.88),(2.24,3.88),(.88,4.66),(2.24,4.66)}{
    \fill[white,draw=VIUDark,line width=.6pt] \p circle (2.4pt);
    \fill[VIUDark] \p circle (1pt);
  }
  \draw[motion] (1.53,4.95)--node[above,font=\fontsize{8}{9.2}\selectfont\bfseries] {$v_L$} (2.65,5.38);
  \node[font=\fontsize{8}{9.2}\selectfont\bfseries,text=VIUDark] at (1.56,3.36)
    {pivotes fijos · $r_i$ constante};

  % (b) Cargo final: el conjunto se traslada y rota sin deformar la formación.
  \draw[motion,dashed] (4.18,3.52).. controls (4.48,3.90) and (4.76,4.15)..
    node[below right,font=\fontsize{8}{9.2}\selectfont\bfseries] {$\Delta p_L$} (5.02,4.34);
  \begin{scope}[rotate around={14:(5.72,4.62)}]
    \pioneericon{5.04}{4.23}{0}{}{.43}
    \pioneericon{6.40}{4.23}{0}{}{.43}
    \pioneericon{5.04}{5.01}{0}{}{.43}
    \pioneericon{6.40}{5.01}{0}{}{.43}
    \draw[rounded corners=.7mm,fill=VIUOrange!18,fill opacity=.72,
      draw=VIUOrange,draw opacity=1,line width=1pt]
      (5.02,4.22) rectangle (6.41,5.02);
    \node[font=\bfseries\fontsize{8}{9.2}\selectfont,rotate=14] at (5.72,4.63) {carga $\mathcal L$};
    \foreach \p in {(5.04,4.23),(6.40,4.23),(5.04,5.01),(6.40,5.01)}{
      \fill[white,draw=VIUDark,line width=.6pt] \p circle (2.4pt);
      \fill[VIUDark] \p circle (1pt);
    }
  \end{scope}
  \draw[SPBlue,->,line width=1pt] (5.10,5.28)
    arc[start angle=160,end angle=62,radius=.42];
  \node[font=\fontsize{8}{9.2}\selectfont\bfseries,text=SPBlue!80!black] at (5.62,5.42)
    {$\Delta\theta_L$};
  \node[font=\fontsize{8}{9.2}\selectfont\bfseries,text=VIUDark] at (5.42,3.36)
    {$r_i$ permanece fijo};

  % (c) Caging inicial: contacto unilateral y carga apoyada en el suelo.
  \draw[SPGreen!70!black,dashed,line width=.75pt,rounded corners=3mm]
    (8.30,3.45) rectangle (10.70,5.16);
  \draw[motion,dashed] (9.48,4.30)--node[above left,font=\fontsize{8}{9.2}\selectfont\bfseries] {$v_L$}
    (10.42,5.16);
  \draw[rounded corners=.7mm,fill=VIUOrange!13,draw=VIUOrange,line width=1pt]
    (8.86,3.90) rectangle (10.16,4.67);
  \node[font=\bfseries\fontsize{8}{9.2}\selectfont] at (9.51,4.29) {carga $\mathcal L$};
  \pioneericon{8.34}{4.28}{0}{}{.46}
  \pioneericon{10.68}{4.28}{180}{}{.46}
  \pioneericon{9.51}{5.13}{-90}{}{.44}
  \pioneericon{9.51}{3.46}{90}{}{.44}
  \draw[unilateral] (8.61,4.28)--(8.82,4.28);
  \draw[unilateral] (10.41,4.28)--(10.20,4.28);
  \draw[unilateral] (9.51,4.88)--(9.51,4.71);
  \draw[unilateral] (9.51,3.72)--(9.51,3.86);
  \node[font=\fontsize{8}{9.2}\selectfont\bfseries,text=SPGreen!48!black] at (9.51,3.27)
    {cerco cerrado · sin unión};

  % (d) Caging final: la carga cambia de pose y los contactos se reubican.
  \draw[motion,dashed] (11.89,3.47).. controls (12.18,3.80) and (12.47,4.05)..
    node[below right,font=\fontsize{8}{9.2}\selectfont\bfseries] {$\Delta p_L$} (12.73,4.27);
  \begin{scope}[rotate around={14:(13.42,4.58)}]
    \draw[SPGreen!70!black,dashed,line width=.75pt,rounded corners=3mm]
      (12.24,3.73) rectangle (14.60,5.42);
    \draw[rounded corners=.7mm,fill=VIUOrange!13,draw=VIUOrange,line width=1pt]
      (12.77,4.19) rectangle (14.07,4.96);
    \node[font=\bfseries\fontsize{8}{9.2}\selectfont,rotate=14] at (13.25,4.38) {$\mathcal L$};

    % Contactos descentrados: sus normales unilaterales generan fuerza y momento.
    \pioneericon{12.29}{4.35}{0}{}{.46}
    \pioneericon{14.55}{4.81}{180}{}{.46}
    \pioneericon{13.15}{5.39}{-90}{}{.44}
    \pioneericon{13.70}{3.76}{90}{}{.44}
    \draw[contactforce,line width=1.55pt]
      (12.54,4.35)--(13.04,4.35);
    \draw[contactforce,line width=.95pt]
      (14.30,4.81)--(14.00,4.81);
    \draw[contactforce,line width=.80pt]
      (13.15,5.15)--(13.15,4.91);
    \draw[contactforce,line width=1.35pt]
      (13.70,4.00)--(13.70,4.45);
    \foreach \p in {(12.78,4.35),(14.06,4.81),(13.15,4.95),(13.70,4.20)}{
      \fill[SPRed] \p circle (1.25pt);
    }

    \draw[resultant] (13.40,4.56)--
      node[below right,font=\fontsize{8}{9.2}\selectfont\bfseries,text=SPBlue] {$F_\Sigma^{E}$} (13.91,4.90);
  \end{scope}
  \draw[VIUOrange,->,line width=1.05pt]
    (13.89,5.15) arc[start angle=18,end angle=145,radius=.32];
  \node[font=\fontsize{8}{9.2}\selectfont\bfseries,text=VIUOrange!85!black] at (13.56,5.29)
    {$\tau_\Sigma^{E}$};
  \node[font=\fontsize{8}{9.2}\selectfont\bfseries,text=SPGreen!48!black] at (13.42,3.27)
    {$F_\Sigma^{E}$ traslada · $\tau_\Sigma^{E}$ gira};
  \end{scope}

  \node[factcard,draw=SPGreen!60!black,fill=SPGreen!7,minimum width=6.30cm,
    minimum height=.84cm] at (3.58,2.88)
    {{\fontsize{8}{9.2}\selectfont\bfseries\color{SPGreen!48!black}CARGO · PIVOTES SOLIDARIOS}\\[-.4mm]
     {\fontsize{8}{9.2}\selectfont $r_i$ fijo; el rumbo del chasis no queda bloqueado}};
  \node[factcard,draw=VIUOrange,fill=VIUOrange!7,minimum width=6.30cm,
    minimum height=.84cm] at (11.42,2.88)
    {{\fontsize{8}{9.2}\selectfont\bfseries\color{VIUOrange!85!black}CAGING · CERCO RECONFIGURABLE}\\[-.4mm]
     {\fontsize{8}{9.2}\selectfont $f_{n,i}^{E}\geq0$: empujes $\rightarrow W_\Sigma^{E}$}};
  \node[circle,fill=white,draw=VIUDark!35,inner sep=1.7pt,font=\bfseries\normalsize]
    at (7.50,2.88) {$\neq$};
  \node[font=\fontsize{8}{9.2}\selectfont\bfseries,text=VIUDark,align=center] at (7.50,1.78)
    {La pose cambia en ambos modos; la restricción física que se conserva es distinta.};

  % Cadena contractual completa: una salida auditable por capa.
  \node[stage,draw=VIUDark,fill=VIUDark!4] (sp1) at (1.06,.50)
    {\textbf{SP1}\\[-.4mm]{\fontsize{6.8}{7.8}\selectfont coalición \textsc{closed}}};
  \node[stage,draw=SPBlue,fill=SPBlue!3] (n1) at (3.64,.50)
    {\textbf{N1 · RUEDAS}\\[-.4mm]{\fontsize{6.8}{7.8}\selectfont $\Xi_{\mathcal C}^{\mathrm{align}}$}};
  \node[stage,draw=VIUOrange,fill=VIUOrange!3,text width=2.04cm] (n2) at (6.22,.50)
    {{\fontsize{7}{8}\selectfont\bfseries\mbox{N2 · CONTACTO}}\\[-.4mm]
     {\fontsize{6.8}{7.8}\selectfont $\widehat{\mathcal W}_{\mathcal C}$}};
  \node[stage,draw=SPPurple,fill=SPPurple!3] (n3) at (8.80,.50)
    {\textbf{N3 · ESTADO}\\[-.4mm]{\fontsize{6.8}{7.8}\selectfont $\hat x_L,\tau_{\mathrm{age}}$}};
  \node[stage,draw=SPGreen,fill=SPGreen!3] (n4) at (11.38,.50)
    {\textbf{N4 · CARGO}\\[-.4mm]{\fontsize{6.8}{7.8}\selectfont EXEC/BRK · HOLD}};
  \node[stage,draw=VIUDark,fill=VIUDark!4] (sp3) at (13.96,.50)
    {\textbf{SP3}\\[-.4mm]{\fontsize{6.8}{7.8}\selectfont $\widehat{\mathcal U}_{\mathcal C}$}};
  \draw[flow] (sp1)--(n1);
  \draw[flow] (n1)--(n2);
  \draw[flow] (n2)--(n3);
  \draw[flow] (n3)--(n4);
  \draw[flow] (n4)--(sp3);
\end{tikzpicture}
}
\captionof{figure}{Modos físicos e interfaz de SP2. Cargo conserva la geometría de los pivotes; caging mantiene un cerco mediante empujes unilaterales. La cadena inferior explicita la salida verificable de SP1, N1--N4 y SP3. Los vectores no están a escala. Elaboración propia.}
\label{fig:sp2-frontera}
\end{center}

\begin{contribucion}{Cadena acreditada de ejecución}
\estadoaporte{N1 demostrado; N2--N4 y Cargo validados empíricamente en planta
reducida; instrumentación y planta Pioneer~3-DX pendientes.}
SP1 entrega la coalición en \textsc{closed}. Las medidas \textsc{observed}
alimentan \textsc{estimated}; el regulador produce \textsc{raw}; y las
comprobaciones de ruedas, soporte, wrench e información deciden el paso a
\textsc{guarded}. Solo N4 alcanza \textsc{executed}, y únicamente en la planta
reducida. Si pierde autoridad, SP2 solicita una reconfiguración a SP1; no ejecuta
otro juego ni una recoalición distribuida. Las herramientas aisladas son
conocidas: la aportación evaluada encadena sus condiciones y mide qué órdenes
dejan de ser ejecutables al incorporar cada capa.
\end{contribucion}

% --------------------------------------------------------------------------
% Modelo mecánico transversal
% --------------------------------------------------------------------------
\clearpage
\leveltag{UNA PLANTA MECÁNICA, DOS LEYES DE CONTACTO}
\section*{Modelo Euler--Lagrange común a Cargo y caging}

Cargo y caging no requieren dos leyes de movimiento incompatibles. Comparten la
dinámica de los AMR y de la carga; difieren en la restricción que genera las
reacciones de contacto. Sea
$q=(q_L,q_1,\ldots,q_N)\in SE(2)^{N+1}$ la configuración del conjunto. En una
carta local, $q$ tiene $3(N+1)$ coordenadas y satisface la ecuación de
Lagrange--d'Alembert
\begin{equation}
M(q)\ddot q+C(q,\dot q)\dot q+D(q)\dot q+g(q)
=B(q)\tau+J_c(q)^\top\lambda+A(q)^\top\gamma_A+d,
\qquad A(q)\dot q=0 .
\label{eq:sp2-unified-el}
\end{equation}
$M$ es definida positiva; $C\dot q$, $D\dot q$ y $g$ reúnen términos
giroscópicos, disipativos y gravitatorios; $B\tau$ contiene los pares de rueda;
$J_c^\top\lambda$, las reacciones de contacto; $A^\top\gamma_A$, las fuerzas que
imponen la rodadura sin deslizamiento lateral; y $d$, las perturbaciones no
modeladas. Para el AMR diferencial $i$, la fila no holónoma es
\begin{equation}
A_i(q_i)=\begin{bmatrix}-\sin\theta_i&\cos\theta_i&0\end{bmatrix},
\qquad -\sin\theta_i\dot x_i+\cos\theta_i\dot y_i=0.
\label{eq:sp2-nonholonomic-row}
\end{equation}
En el plano horizontal, la gravedad no entra en las tres coordenadas de
$q_L$; su efecto vertical reaparece en el reparto de soporte de
\eqref{eq:vertical-support}.

\paragraph{Cargo: contacto bilateral.}
El pivote soportado por el robot $i$ conserva el desplazamiento $r_i$ respecto
de la carga:
\begin{equation}
\phi_i^{C}(q)=p_i-p_L-R(\theta_L)r_i=0,
\qquad J_C(q)\dot q=0,
\qquad J_C=\frac{\partial\phi^{C}}{\partial q}.
\label{eq:sp2-cargo-bilateral}
\end{equation}
La reacción bilateral puede mantener la unión en ambas direcciones dentro de
los límites del acoplamiento. Al proyectar \eqref{eq:sp2-unified-el} sobre la
carga se obtiene el modelo planar usado por la guarda mecánica,
\begin{equation}
M_L\ddot q_L+D_L\dot q_L
=T(\theta_L)G_{\mathcal C}\mathbf f+W_{\mathrm{ext}},
\label{eq:sp2-payload-el}
\end{equation}
donde $G_{\mathcal C}$ transforma las reacciones individuales $\mathbf f$ en
wrench de la carga. La matriz puede permanecer fija en el marco solidario
mientras no cambien los contactos.

\paragraph{Caging: contacto unilateral.}
En empuje no se impone \eqref{eq:sp2-cargo-bilateral}. Si $g_i^{E}(q)$ es la
separación normal, las reacciones cumplen
\begin{equation}
0\leq f_{n,i}^{E}\ \perp\ g_i^{E}(q)\geq0,
\qquad |f_{t,i}^{E}|\leq\mu_i f_{n,i}^{E},
\qquad
\mathbf f_i^{E}=f_{n,i}^{E}n_i^{E}+f_{t,i}^{E}t_i^{E}.
\label{eq:sp2-caging-complementarity}
\end{equation}
La primera relación expresa complementariedad: una separación positiva anula
la fuerza normal; una fuerza normal positiva exige contacto. Para el conjunto
activo $\mathcal A(t)$, la carga obedece
\begin{equation}
M_L\ddot q_L+D_L\dot q_L
=\sum_{i\in\mathcal A(t)}G_i^{E}(q_L)\mathbf f_i^{E}+W_{\mathrm{ext}}.
\label{eq:sp2-caging-el}
\end{equation}
El cambio de $\mathcal A(t)$ convierte el modelo en híbrido. Estas ecuaciones
determinan cómo se mueve la carga cuando existen fuerzas; no prueban que el
objeto esté encerrado. Esa propiedad exige que la componente libre que contiene
$q_L$ en $SE(2)$ no alcance una configuración de escape
\parencite{bicchi1995closure,pereira2004caging}. El verificador disponible al
final de SP2 fija $\theta_L$ y explora solo $(x_L,y_L)$; por tanto, no acredita
caging completo ni que tres robots basten. La formulación tampoco incorpora una
ley de impacto o restitución: describe fases persistentes de contacto y su
conmutación.

\section*{Caso nominal: carga cuadrada sostenida por tres AMR}

El siguiente caso cierra la cadena fuerza--wrench--movimiento sin convertirla en
resultado experimental. Considérese una lámina cuadrada uniforme de lado $a$ y
masa $m_L$. Su inercia polar es $J_L=m_La^2/6$ y
$M_L=\operatorname{diag}(m_L,m_L,J_L)$. Los apoyos
\begin{equation}
r_1=\begin{bmatrix}a/2\\0\end{bmatrix},\qquad
r_2=\begin{bmatrix}-a/2\\a/2\end{bmatrix},\qquad
r_3=\begin{bmatrix}-a/2\\-a/2\end{bmatrix}
\label{eq:sp2-square-supports}
\end{equation}
contienen el centro de masa en su triángulo. El equilibrio vertical da
\begin{equation}
n_1=\frac{m_Lg_0}{2},\qquad
n_2=n_3=\frac{m_Lg_0}{4}.
\label{eq:sp2-square-normal-loads}
\end{equation}
Para $m_L=18$ kg y $g_0=9{,}81$ m/s$^2$, las reacciones son
$88{,}29$, $44{,}15$ y $44{,}15$ N. La geometría, no el número de robots,
determina el reparto.

Con $\mathbf f=(f_{x,1},f_{y,1},\ldots,f_{x,3},f_{y,3})$, la matriz de agarre
\parencite{nakamura1989dynamics,yoshikawa1993coordinated} es
\begin{equation}
G_3=\begin{bmatrix}
1&0&1&0&1&0\\
0&1&0&1&0&1\\
0&a/2&-a/2&-a/2&a/2&-a/2
\end{bmatrix},
\qquad W_C^{L}=G_3\mathbf f .
\label{eq:sp2-square-grasp}
\end{equation}
Así se separan dos decisiones: el servocontrol calcula el wrench requerido y el
repartidor busca fuerzas admisibles que lo realicen.

Sea $T_f>0$ la duración y $s=t/T_f\in[0,1]$. El polinomio
$h(s)=10s^3-15s^4+6s^5$ anula velocidad y aceleración en ambos extremos. Una
referencia suave entre $q_0$ y $q_f$ es
\begin{equation}
q_d(t)=q_0+h(t/T_f)
\begin{bmatrix}
x_f-x_0\\ y_f-y_0\\
\operatorname{wrap}(\theta_f-\theta_0)
\end{bmatrix}.
\label{eq:sp2-quintic-reference}
\end{equation}
Defínanse $e=q_L-q_d$, con el error angular envuelto en $(-\pi,\pi)$, y
$\dot e=\dot q_L-\dot q_d$. Para el modelo nominal de
\eqref{eq:sp2-payload-el}, el wrench de par calculado es
\begin{equation}
W_d^0=M_L\ddot q_d+D_L\dot q_L-K_D\dot e-K_Pe,
\qquad G_3\mathbf f=T(\theta_L)^\top W_d^0.
\label{eq:sp2-computed-wrench}
\end{equation}

\Needspace{17\baselineskip}
\begin{sp2prop}[Estabilidad del Cargo nominal]
Supónganse contactos bilaterales fijos, $M_L\succ0$ constante,
$K_P\succ0$, $K_D\succ0$, ausencia de perturbaciones y saturaciones, y
realización exacta de \eqref{eq:sp2-computed-wrench}. Mientras el error angular
permanezca en la carta $(-\pi,\pi)$, el equilibrio $(e,\dot e)=(0,0)$ es
asintóticamente estable.
\end{sp2prop}
\begin{proof}
Sustituir \eqref{eq:sp2-computed-wrench} en
\eqref{eq:sp2-payload-el} produce
$M_L\ddot e+K_D\dot e+K_Pe=0$. Para
\begin{equation}
V(e,\dot e)=\tfrac12\dot e^\top M_L\dot e+
\tfrac12e^\top K_Pe
\label{eq:sp2-cargo-lyapunov}
\end{equation}
se obtiene $\dot V=-\dot e^\top K_D\dot e\leq0$. El mayor conjunto invariante
contenido en $\{\dot e=0\}$ satisface $K_Pe=0$; por tanto, $e=0$, y el principio
de invariancia concluye la convergencia.
\end{proof}

El resultado no se transfiere al caging, a cambios de contacto ni a la planta
Pioneer~3-DX sin identificar. Tampoco añade un controlador al benchmark N4:
formaliza el caso nominal que la guarda intenta aproximar mediante órdenes de
aceleración, soporte secuencial y reparto de wrench.

% --------------------------------------------------------------------------
% Alcance y comparación metodológica
% --------------------------------------------------------------------------
\leveltag{EJECUCIÓN CAUSAL BAJO RESTRICCIONES ACOPLADAS}
\section*{Objetivo de control y límites de validez}

En Cargo, la referencia de pose se convierte en ejecución nominal, frenado
activo, mantenimiento de una carga ya detenida o declaración de pérdida de
autoridad. La variable $m$ denota el índice temporal; $k$ se reserva para las
cargas, de acuerdo con la notación general del TFM.

\paragraph{Ejecución causal y reconfiguración.}
Sean $x_{L,m}=(q_{L,m},\xi_{L,m})$, la coalición vigente $\mathcal C_m$, la
referencia $q_L^\star$ y la historia local $\mathcal I_{i,m}$. En el modelo
reducido, las cinco leyes no predictivas se calculan en un AMR de referencia $\ell$
(el primero de la coalición) a partir de su última estimación. El MPC recibe el
estado global y actúa solo como techo central. La actualización ejecutada es
\begin{equation}
\begin{aligned}
a_{L,m}^0&=
\begin{cases}
\pi_\ell(\mathcal I_{\ell,m},q_L^\star), & \text{cinco leyes locales},\\
\pi_{\mathrm{MPC}}(x_{L,m},q_L^\star), & \text{MPC central reducido},
\end{cases}\\
x_{L,m+1}&=\mathsf F_{\Delta t}(x_{L,m},a_{L,m},d_m),\quad
\xi_{L,m+1}\in\Xi_{\mathcal C_m}^{\mathrm{align}},\\
W_d(x_{L,m},a_{L,m})&\in\widehat{\mathcal W}_{\mathcal C_m},\quad
g_{\mathrm{info},m}=1 .
\end{aligned}
\label{eq:sp2-execution-problem}
\end{equation}
El objetivo operativo es hallar el primer instante $M$ que alcanza el terminal y
mantiene las restricciones de \eqref{eq:sp2-execution-problem} desde $m=0$
hasta $M-1$. Las políticas ensayadas no optimizan $M\Delta t$ de forma global.
Para las leyes locales,
$g_{\mathrm{info},m}=\mathbf 1[\tau_{\mathrm{age},\ell,m}\leq\bar\tau]$; en
MPC vale uno por definición. En los escenarios con obstáculo, la guarda muestreada
propone $\widetilde a_{L,m}^0=\mathsf B_m(x_{L,m},a_{L,m}^0)$; en los demás,
$\widetilde a_{L,m}^0=a_{L,m}^0$. La orden ejecutada la determina el modo
$\sigma_m\in\{\mathrm{EXECUTE},\mathrm{BRAKE},\mathrm{HOLD},
\mathrm{UNCONTROLLED}\}$:
\begin{equation}
a_{L,m}=\begin{cases}
\alpha_m \widetilde a_{L,m}^0, & \mathrm{EXECUTE},\quad \alpha_m>0,\\
\beta_m a_{L,m}^{\mathrm{br}}, & \mathrm{BRAKE},\quad
a_{L,m}^{\mathrm{br}}=-K_v\xi_{L,m},\ \beta_m>0,\\
0, & \mathrm{HOLD},\quad \|\xi_{L,m}\|\leq\varepsilon_v,\\
0, & \mathrm{UNCONTROLLED}\quad\text{(sin orden certificable)}.
\end{cases}
\label{eq:sp2-operational-modes}
\end{equation}
En \eqref{eq:sp2-operational-modes}, cada candidato nominal o de frenado se
vuelve a comprobar conjuntamente frente
a ruedas, soporte, $\widehat{\mathcal W}_{\mathcal C_m}$, edad de información y
guarda de obstáculo. En \textsc{uncontrolled}, el cero representa ausencia de
autoridad comandada en la planta reducida, no una parada física garantizada; se
registra el fallo y se solicita reconfiguración.
El estado debe permanecer durante $0{,}80$ s en
\begin{equation}
\mathcal E=\left\{x_L:
\begin{aligned}
\|p_L-p_L^\star\|_2&\leq0{,}10\ \mathrm m,
&|\operatorname{wrap}(\theta_L-\theta_L^\star)|&\leq0{,}10\ \mathrm{rad},\\
\|v_L\|_2&\leq0{,}12\ \mathrm{m/s},
&|\omega_L|&\leq0{,}10\ \mathrm{rad/s}
\end{aligned}\right\}.
\label{eq:sp2-terminal-set}
\end{equation}
La maniobra termina cuando el estado permanece en
\eqref{eq:sp2-terminal-set} durante el intervalo prescrito. El mismo criterio se
aplica a todas las aproximaciones causales y a los mismos mundos. Los resultados
describen el modelo reducido; la optimalidad global y la
estabilidad del sistema híbrido completo quedan fuera de esta comparación.

La Tabla~\ref{tab:sp2-methods} sitúa la comparación. Ninguno de los métodos que
resume cubre a la vez movilidad diferencial, contacto, información local y
control de la carga; su última columna identifica la restricción ausente. Se
denomina \emph{wrench} planar al vector conjunto de fuerza
longitudinal, fuerza transversal y momento alrededor de la carga,
$W=(F_x,F_y,\tau_z)$.
\begin{table}[H]
\centering
\fontsize{8.0}{8.55}\selectfont
\setlength{\tabcolsep}{3.2pt}
\renewcommand{\arraystretch}{0.88}
\caption{Métodos, papel experimental y restricción que no resuelven por sí solos.}
\label{tab:sp2-methods}
\begin{tabularx}{\textwidth}{p{2.05cm} p{2.65cm} p{3.25cm} X}
\toprule
Método & Papel en SP2 & Información y arquitectura & Garantía o límite pertinente \\
\midrule
Líder--seguidor \parencite{desai2001formation} & Referencia distribuida & Estado del líder sobre un árbol & Trata movilidad no holónoma; no certifica soporte ni wrench \\
Estructura virtual \parencite{ren2004virtual} & Estimación/consenso & Referencia común en grafo vecinal & El error depende de conectividad y canal; no representa contacto \\
PD/PID & Baseline causal de pose & Última estimación del AMR de referencia & Sin restricciones explícitas; PID añade integral acotada \\
LQR/LQI \parencite{kalman1960optimal} & Sustituto lineal & Doble integrador local; LQI añade integral & Optimalidad solo para el modelo y pesos fijados \\
Rigidez--amortiguamiento \parencite{ortega2002passivity} & Baseline energético & Error de pose y velocidad local & No es una realización port--Hamiltoniana ni prueba pasividad \\
MPC \parencite{mayne2000mpc} & Techo central & Estado global; horizonte 8 & Restringe aceleración, no contactos; no es DMPC/NMPC \\
Conmutación \parencite{lee2025switching} & Contraste de literatura & MILP--QP central entre eventos & Sus garantías no se transfieren al contacto de Cargo \\
Guarda SP2 & Capa transversal & N1 + N2 + edad de N3 & Certifica el candidato muestreado; no la trayectoria física continua \\
\bottomrule
\end{tabularx}
\end{table}

La equivalencia cinemática, la factibilidad secuencial de soporte y wrench, la
calidad informativa bajo canal local y la comparación pareada de control se
establecen bajo supuestos distintos. En
particular, la restauración de S4 es exógena y el MPC recibe estado global.

% --------------------------------------------------------------------------
% Página 3 · N1 formulación
% --------------------------------------------------------------------------
\leveltag{COMPATIBILIDAD ENTRE CUERPO RÍGIDO Y RUEDAS}
\section*{Condiciones para el movimiento de cuerpo rígido}

El modelo cinemático omite masa y fricción. Supone seguimiento exacto y
acoplamientos sin holgura. El pivote de cada AMR se sitúa en el punto medio de su eje. El
acoplamiento fija ese punto respecto a la carga y permite el giro propio del
AMR; una unión que también bloquee el rumbo requiere otro modelo. Si $r_i$ es
la posición del pivote en el marco de la carga,
$J=\left[\begin{smallmatrix}0&-1\\1&0\end{smallmatrix}\right]$ y
$R(\theta_L)$ la rotación plana, entonces
\begin{align}
v_i^{\mathrm{piv}}&=v+\omega J R(\theta_L)r_i,
\label{eq:pivot-velocity}\\
a_i^{\mathrm{piv}}&=\dot v+\dot\omega J R(\theta_L)r_i
-\omega^2R(\theta_L)r_i.
\label{eq:pivot-acceleration}
\end{align}
Las expresiones \eqref{eq:pivot-velocity}--\eqref{eq:pivot-acceleration} fijan
la cinemática del pivote. Sea
$\mathbf e_i=(\cos\theta_i,\sin\theta_i)$. La restricción no holónoma exige
una velocidad longitudinal con signo $v_{i,\parallel}$ tal que
\begin{equation}
v_i^{\mathrm{piv}}=v_{i,\parallel}\mathbf e_i.
\label{eq:heading-alignment}
\end{equation}
La rama $v_{i,\parallel}>0$ corresponde a avance y
$v_{i,\parallel}<0$ a marcha atrás. El rumbo
admisible puede ser la dirección de $v_i^{\mathrm{piv}}$ o esa misma dirección
más $\pi$. Cuando $\|v_i^{\mathrm{piv}}\|>0$, ambas ramas comparten la tasa
\begin{equation}
\dot\theta_i^\star=
\frac{v_{i,x}^{\mathrm{piv}}a_{i,y}^{\mathrm{piv}}-
v_{i,y}^{\mathrm{piv}}a_{i,x}^{\mathrm{piv}}}
{\|v_i^{\mathrm{piv}}\|^2}.
\label{eq:heading-rate}
\end{equation}
Si $r_i^w$, $\ell_i^w$ y $\bar\omega_{w,i}$ denotan el radio de rueda, la
semivía y el límite de velocidad angular, las ruedas deben satisfacer
\begin{equation}
\dot\varphi_{i,L/R}=\frac{1}{r_i^w}
\left(v_{i,\parallel}\mp\ell_i^w\dot\theta_i^\star\right),
\qquad |\dot\varphi_{i,L/R}|\leq\bar\omega_{w,i}.
\label{eq:wheel-speeds}
\end{equation}

La equivalencia anterior es exclusivamente cinemática. La dinámica de cada
Pioneer~3-DX está contenida en \eqref{eq:sp2-unified-el}; para evaluarla en
hardware aún deben identificarse inercia, disipación, pares de rueda y reacción
del acoplamiento. La campaña actual no proporciona esos parámetros.

El margen de viabilidad cinemática de la orden se escribe como
\begin{equation}
m_{\mathrm{kin}}=
1-\max_{i\in\mathcal C}\max_{s\in\{L,R\}}
\frac{|\dot\varphi_{i,s}|}{\bar\omega_{w,i}},
\label{eq:kinematic-margin}
\end{equation}
donde $\bar\omega_{w,i}$ es el límite de rueda del AMR $i$. En
\eqref{eq:kinematic-margin}, $m_{\mathrm{kin}}>0$ indica holgura,
$m_{\mathrm{kin}}=0$ corresponde a
saturación y $m_{\mathrm{kin}}<0$ rechaza la orden. La definición permanece
finita cuando una rueda está detenida y comparte el umbral cero con el margen
mecánico definido más adelante. El registro conserva el AMR y la rueda limitantes.

\begin{sp2prop}[Equivalencia cinemática]
Considérese un intervalo en el que los pivotes tienen velocidad distinta de cero
y la velocidad de la carga es continuamente diferenciable. Bajo los supuestos
anteriores, supóngase además que cada rumbo inicial satisface
\eqref{eq:heading-alignment}. La coalición reproduce $\xi_L$ si y solo si cada
AMR admite una rama continua $(v_{i,\parallel},\theta_i)$ que cumple desde
\eqref{eq:pivot-velocity} hasta \eqref{eq:wheel-speeds}.
\end{sp2prop}
\begin{proof}
Derivar $p_i=p_L+R(\theta_L)r_i$ da
\eqref{eq:pivot-velocity}. La cinemática diferencial obliga a
\eqref{eq:heading-alignment} y transforma
$(v_{i,\parallel},\dot\theta_i)$ en las dos
velocidades de rueda de \eqref{eq:wheel-speeds}. Recíprocamente, esas igualdades
realizan la derivada de $p_i$ sin componente lateral. Como la posición inicial
del pivote coincide con el anclaje y la rama de rumbo es continua, la relación
rígida se conserva en el intervalo.
\end{proof}

\begin{sp2lema}[Ventana de validez]
El conjunto realizable es
$\Xi_{\mathcal C}^{\mathrm{align}}=
\bigcap_{i\in\mathcal C}\Xi_i^{\mathrm{align}}$. En un giro puro alrededor de
un centro instantáneo, con $d_i^{\mathrm{ICR}}$ igual a la distancia del pivote
a ese centro,
\begin{equation}
|\omega|\leq
\min_{i\in\mathcal C}
\frac{r_i^w\bar\omega_{w,i}}{d_i^{\mathrm{ICR}}+\ell_i^w}.
\label{eq:pure-rotation-bound}
\end{equation}
\end{sp2lema}
La rueda exterior del AMR con menor cociente fija la cota. Un pivote inmóvil
vuelve singular \eqref{eq:heading-rate} y deja el caso fuera de la proposición.
También quedan fuera la aceleración de rueda, el par motor y la estabilidad del
contacto.

La cadena necesaria queda fijada por las ecuaciones anteriores: geometría rígida
en el anclaje, velocidad sin componente lateral y conversión a las dos ruedas.
El pivote es una idealización cinemática, no un acoplamiento mecánico adicional.


% --------------------------------------------------------------------------
% Página 4 · N1 experimento
% --------------------------------------------------------------------------
\leveltag{ENVOLVENTE DE MOVILIDAD DE LA FORMACIÓN}
\section*{Barrido de radio y velocidad angular}

La formación cardinal se evaluó con radios entre $0{,}20$ y $1{,}00$ m,
velocidades longitudinales entre $0$ y $0{,}30$ m/s y 61 valores de $\omega$
hasta $1{,}50$ rad/s. El radio de rueda fue $0{,}06$ m y la vía $0{,}32$ m. Para
cada una de 30 semillas, los límites nominales $12$, $10$, $8$ y $6$ rad/s se
perturbaron de forma independiente en un intervalo de $\pm18\%$. El comparador
básico solo compara la rapidez del pivote con
$r_i^w\bar\omega_{w,i}$; la comprobación
completa calcula la tasa de rumbo y ambas ruedas.

Los paneles A, B y D de la Figura~\ref{fig:sp2-master-n1-n2} cuantifican cómo el
radio y la velocidad longitudinal reducen el margen angular y cuántas órdenes
acepta por error el comparador básico.

\textbf{Resultado N1.} Comprobar solo la rapidez de cada pivote sobreestima la
movilidad de la coalición. El barrido contiene \SPTwoReadyNOneRows{} órdenes y
el comparador básico acepta
\SPTwoReadyNOneFalse{} que la comprobación conjunta rechaza: el 12,4~\% de las
órdenes ensayadas. El análisis identifica el AMR y la rueda limitantes en cada
rechazo. La fracción factible disminuye de forma monótona entre el menor y el
mayor radio ensayados, en consonancia con \eqref{eq:pure-rotation-bound}. Las
proporciones pertenecen a la cuadrícula fijada y no definen una ley asintótica
para una población de robots.

\leveltag{RETIRADA DE UN ROBOT Y MONOTONÍAS OPUESTAS}
\section*{Efecto de retirar un AMR}

Si el AMR $j$ abandona la carga, la condición cinemática pasa a ser
\begin{equation}
\xi_L\in\Xi_{\mathcal C\setminus\{j\}}^{\mathrm{align}}
=\bigcap_{i\in\mathcal C\setminus\{j\}}\Xi_i^{\mathrm{align}}.
\label{eq:n1-removal}
\end{equation}
La intersección de \eqref{eq:n1-removal} crece al retirar una restricción cinemática. La capacidad
mecánica, en cambio, evoluciona en la dirección opuesta.

\begin{sp2prop}[Monotonicidades opuestas]
Sea $j\in\mathcal C$ y supóngase que la reacción nula pertenece al conjunto local
admisible de cada contacto. Si $\mathcal W_{\mathcal C}$ denota el conjunto
físico completo, unión sobre todos los repartos normales factibles, entonces
\begin{equation}
\Xi_{\mathcal C\setminus\{j\}}^{\mathrm{align}}
\supseteq\Xi_{\mathcal C}^{\mathrm{align}},
\qquad
\mathcal W_{\mathcal C\setminus\{j\}}\subseteq\mathcal W_{\mathcal C}.
\label{eq:opposite-monotonicity}
\end{equation}
Por tanto, el conjunto de órdenes cinemática y mecánicamente admisibles no es,
en general, monótono con el número de robots.
\end{sp2prop}
\begin{proof}
La primera inclusión resulta de eliminar $\Xi_j^{\mathrm{align}}$ de una
intersección. Para
la segunda, cualquier reparto de los AMR restantes también es un reparto de
$\mathcal C$ si se asigna reacción nula a $j$. La inclusión no se transfiere al
oráculo secuencial de contacto, que vuelve a seleccionar un reparto vertical al cambiar
la coalición. No existe una inclusión uniforme
para la condición conjunta. Si $j$ limita las ruedas pero su reacción óptima es
nula, retirarlo amplía el conjunto de órdenes admisibles. Si $j$ no limita las ruedas pero su
reacción es imprescindible para equilibrar el peso o el momento, retirarlo lo
reduce. Estos dos casos construyen las dos direcciones posibles.
\end{proof}

La retirada exige repetir ambas comprobaciones. Con límites de rueda
heterogéneos tampoco es válido sustituir la intersección por una media: basta que
un AMR exceda su cota para rechazar la orden común. El enunciado describe el
estado posterior a la retirada; no modela el transitorio de desacoplamiento.

% --------------------------------------------------------------------------
% Página 5 · N2 formulación
% --------------------------------------------------------------------------
\leveltag{SOPORTE, FRICCIÓN Y EQUILIBRIO DE LA CARGA}
\section*{Soporte vertical y equilibrio plano}

La capa mecánica enlaza la aceleración solicitada con ese \emph{wrench} de
contacto. En el marco inercial se calcula
\begin{equation}
W_d^0=M_La_L^0+D_L\xi_L-\widehat W_{\mathrm{ext}}^0,
\qquad
W_d=T(\theta_L)^\top W_d^0,
\label{eq:load-dynamics}
\end{equation}
donde $T(\theta_L)=\operatorname{diag}(R(\theta_L),1)$ expresa el wrench en el
marco de la carga. El signo negativo evita contar dos veces una perturbación ya
incluida en la dinámica. La implementación registra $W_d^0$ y $W_d$ antes de
llamar al repartidor.

El modelo 2.5D mantiene la carga horizontal y resuelve primero el soporte
vertical. Si $n_i$ es la reacción normal del AMR $i$, $r_i=(x_i,y_i)$ su
posición y $(x_c,y_c)$ la proyección del centro de masa, el equilibrio
cuasiestático exige
\begin{equation}
\sum_{i\in\mathcal C} n_i=m_Lg,
\qquad
\sum_i(x_i-x_c)n_i=0,
\qquad
\sum_i(y_i-y_c)n_i=0,
\qquad 0\leq n_i\leq\bar n_i.
\label{eq:vertical-support}
\end{equation}
Siguiendo la matriz de agarre planar usada en manipulación cooperativa
\parencite{nakamura1989dynamics,yoshikawa1993coordinated}, una reacción
$\mathbf t_i=(t_{i,x},t_{i,y})$, expresada en el marco de la carga, contribuye a
\begin{equation}
W_{\mathcal C}=G_{\mathcal C}\mathbf t,
\qquad
G_{\mathcal C}=\begin{bmatrix}
1&0&\cdots&1&0\\
0&1&\cdots&0&1\\
-r_{1,y}&r_{1,x}&\cdots&-r_{n,y}&r_{n,x}
\end{bmatrix}.
\label{eq:grasp-matrix}
\end{equation}
La matriz \eqref{eq:grasp-matrix} conserva fuerza y momento. El oráculo ejecuta dos etapas. Primero obtiene $n^\star$ resolviendo
\begin{equation}
\min_n\ \sum_i\left(\frac{n_i-n_i^{\mathrm{nom}}}{m_Lg}\right)^2
\quad\text{sujeto a \eqref{eq:vertical-support}},
\label{eq:vertical-support-program}
\end{equation}
con $n_i^{\mathrm{nom}}=m_Lg/|\mathcal C|$ recortado a las cotas del contacto.
El problema \eqref{eq:vertical-support-program} es un QP convexo con
restricciones lineales. La implementación reutiliza SLSQP para conservar una
única interfaz de restricciones acotadas, sin atribuirle una ventaja algorítmica.
La solución solo se acepta si supera el control de residual siguiente. Para no
mezclar N y N\,m, sea
$L_{\mathcal C}=\max\{0{,}1\ \mathrm m,\max_i\|r_i-r_c\|_2\}$ y defínase
\begin{equation}
r_{\mathrm{sup}}=\left\|\left(
\frac{\sum_i n_i-m_Lg}{m_Lg},
\frac{\sum_i(y_i-y_c)n_i}{m_LgL_{\mathcal C}},
\frac{\sum_i(x_i-x_c)n_i}{m_LgL_{\mathcal C}}
\right)\right\|_\infty .
\label{eq:support-normalized-residual}
\end{equation}
El reparto se acepta si el residual de \eqref{eq:support-normalized-residual}
satisface $r_{\mathrm{sup}}\leq10^{-7}$.

Después se fijan esas reacciones. Aquí $\mu_{L,i}$ corresponde al contacto
carga--plataforma. En un modelo identificado, $\alpha_i\in(0,1]$ sería la
fracción de la reacción rueda--suelo soportada por el eje motriz y
$N_{g,i}=\alpha_i(m_i g+n_i^\star)$. El experimento no mide $\alpha_i$, la
adherencia lateral ni el par del Pioneer: fija directamente $\bar t_i$ desde
la configuración experimental. Una extensión
anisótropa distinguiría, para el rumbo $\mathbf e_i$, las cotas
$|\mathbf e_i^\top\mathbf t_i|\leq F_{\mathrm{drive},i}$ y
$|\mathbf e_i^{\perp\top}\mathbf t_i|\leq F_{\mathrm{lat},i}$. El LP ejecutado usa en
cambio un radio isotrópico inscrito; es un sustituto 2.5D, no una identificación
de la envolvente de tracción del Pioneer 3-DX.
La capacidad tangencial inscrita del contacto $i$ es
\begin{equation}
\ell_i=\cos(\pi/16)\min(\mu_{L,i} n_i^\star,\bar t_i),\qquad
\mathcal T_i=\{\mathbf t_i:d_h^\top\mathbf t_i\leq\ell_i, h=1,\ldots,16\},
\label{eq:local-force-set}
\end{equation}
donde, en \eqref{eq:local-force-set}, $d_h$ son normales unitarias. Si $\mathcal N_{\mathcal C}$ reúne todos
los repartos normales que satisfacen \eqref{eq:vertical-support}, el conjunto
físico completo y la sección que usa el oráculo son, respectivamente,
\begin{equation}
\mathcal W_{\mathcal C}=\bigcup_{n\in\mathcal N_{\mathcal C}}
\{G_{\mathcal C}\mathbf t:\mathbf t_i\in\mathcal T_i(n_i)\},\qquad
\widehat{\mathcal W}_{\mathcal C}=\{G_{\mathcal C}\mathbf t:
\mathbf t_i\in\mathcal T_i(n_i^\star)\}.
\label{eq:physical-sequential-wrench-sets}
\end{equation}
La Proposición de monotonías opuestas se refiere a $\mathcal W_{\mathcal C}$ de
\eqref{eq:physical-sequential-wrench-sets}.
El oráculo secuencial vuelve a calcular $n^\star$ al cambiar la coalición y no
hereda necesariamente esa inclusión. El objetivo vertical, por su parte,
aproxima un reparto uniforme en vez de maximizar la capacidad tangencial: otro reparto normal
factible podría admitir un wrench que la sección
$\widehat{\mathcal W}_{\mathcal C}$ rechaza.

El primer LP minimiza la utilización con balance exacto:
\begin{equation}
\begin{aligned}
\rho^\star=\min_{\mathbf t,\rho}\quad &\rho,\\
\text{sujeto a}\quad &G_{\mathcal C}\mathbf t=W_d,\\
&d_h^\top\mathbf t_i\leq\rho\ell_i,\quad \rho\geq0.
\end{aligned}
\label{eq:wrench-allocation}
\end{equation}
La orden se acepta si la etapa vertical es factible, el balance normalizado es
menor que $\epsilon_W$ y $\rho^\star\leq1+\epsilon_\rho$. Si se rechaza, un
segundo LP, con $\rho\leq1$, minimiza
$\|Q_W(s^++s^-)\|_1$ sujeto a
$G_{\mathcal C}\mathbf t+s^+-s^-=W_d$. Esta segunda solución solo diagnostica
el wrench físicamente alcanzable; no puede convertir el rechazo en aceptación.
Así desaparece la dependencia de un peso entre utilización y holgura. Se guarda
\[
m_{\mathrm{wrench}}=1-\rho^\star,
\quad P_{\ker G_{\mathcal C}}=I-G_{\mathcal C}^{\dagger}G_{\mathcal C},
\quad \rho_{\mathrm{int}}=
\frac{\|P_{\ker G_{\mathcal C}}\mathbf t\|_2}{\|\mathbf t\|_2+\varepsilon}.
\]
Las demandas aceptadas pertenecen a $\widehat{\mathcal W}_{\mathcal C}$; las rechazadas
conservan su utilización o residual para explicar el fallo. El procedimiento usa
información global y actúa como oráculo mecánico
\parencite{nakamura1989dynamics,yoshikawa1993coordinated}.

% --------------------------------------------------------------------------
% Página 6 · N2 experimento
% --------------------------------------------------------------------------
\leveltag{CAPACIDAD DE WRENCH BAJO DEMANDA Y DISPERSIÓN}
\section*{Capacidad escalar, LP bilateral y contacto 2.5D}

El experimento compara el oráculo de \eqref{eq:wrench-allocation} con dos
relajaciones. La primera descarta geometría y acepta cuando la
capacidad escalar supera la fuerza solicitada. La segunda usa el LP bilateral
\begin{equation}
\min_{\mathbf t,\eta^{N2}}\ \eta^{N2}
\quad\text{sujeto a}\quad
G_{\mathcal C}\mathbf t=W_d,\quad
|t_{i,x}|,|t_{i,y}|\leq\eta^{N2}\bar t_i,\quad \eta^{N2}\geq0.
\label{eq:wrench-lp-pilot}
\end{equation}
El LP \eqref{eq:wrench-lp-pilot} comprueba el momento, pero no el peso ni la fricción. El tercer modelo
resuelve soporte, polígono de Coulomb y límites de tracción heterogéneos. Las
posiciones rectangulares reciben ruido gaussiano de $0{,}025$ m. Se fijaron 30
semillas por celda y los factores
\begin{equation}
d\in\{0{,}20,0{,}40,0{,}60,0{,}80\}\ \mathrm{m},
\pi^{N2}\in\{0{,}2,0{,}4,0{,}6,0{,}8,1{,}0,1{,}2\},
\label{eq:n2-pilot-factors}
\end{equation}
Con los factores de \eqref{eq:n2-pilot-factors}, la demanda es
\begin{equation}
W_d(\pi^{N2})=(40\pi^{N2},\ 10\pi^{N2},\ 12\pi^{N2}).
\label{eq:n2-pilot-demand}
\end{equation}
La escala $\pi^{N2}$ es adimensional; las dos primeras componentes se expresan
en N y la tercera, en N\,m. La dispersión $d$ controla el brazo disponible para
producir momento.
La masa de la carga es $18$ kg. Las 30 semillas describen variación geométrica
dentro de la cuadrícula ensayada; no reproducen incertidumbre de contacto medida.

La Figura~\ref{fig:sp2-master-n1-n2} reúne la envolvente N1, sus falsos
positivos, los tres modelos N2 y la lectura conjunta de ambos barridos. El panel
C mantiene la misma cuadrícula de demanda y geometría al añadir soporte y
fricción.
\resultfigure[0.55\textheight]{figures/generated/fig-sp2-master-n1-n2.pdf}{De coalición seleccionada a coalición ejecutable. A: envolvente de ruedas. B: falsa aceptación del criterio simple. C: escalera de modelos mecánicos. D: efecto opuesto de la escala geométrica en los dos barridos; no es una estimación causal común.}{fig:sp2-master-n1-n2}

\textbf{Resultado N2.} Tener fuerza total suficiente no garantiza producir la
fuerza y el momento solicitados. La capacidad escalar y el LP bilateral aceptan
las \SPTwoReadyNTwoRows{}
demandas. El oráculo 2.5D acepta \SPTwoReadyNTwoFeasible{} y rechaza
\SPTwoReadyNTwoFalse{}, el \spdecimal{\SPTwoReadyNTwoRejectRate}~\% del
barrido, que ambos modelos simplificados habían aceptado. En esta cuadrícula,
el LP bilateral coincide con la capacidad escalar en todos los casos. El
contraejemplo de momento prueba que no son equivalentes en general, aunque la
campaña no observó esa diferencia. Los rechazos del modelo 2.5D se
concentran en el menor brazo ensayado, donde el momento exige fuerzas que
superan la adherencia de los contactos débiles. El residual y la utilización
quedan almacenados por ejecución; los fallos permanecen en el denominador.

\leveltag{FRONTERA MECÁNICA Y SUPUESTOS DE CONTACTO}
\section*{Necesidad y suficiencia de las comprobaciones}

El barrido fija la orientación de los contactos, el centro de masa y la dirección
de la fuerza. Esos paneles cubren, en consecuencia, una sola
familia de wrenches y no cualquier maniobra.

\paragraph{Contraejemplo de capacidad escalar.}
Si $\bar t_i$ acota la norma de cada reacción, la desigualdad triangular hace
necesaria la condición $\sum_i\bar t_i\geq\|F_d\|$. No es suficiente para
producir $W_d=(F_d,\tau_d)$: si todos los contactos coinciden con el centro,
$r_i\times\mathbf t_i=0$ y cualquier $\tau_d\neq0$ resulta inviable, por grande
que sea la capacidad escalar. Las relajaciones de
\eqref{eq:wrench-allocation} contienen el conjunto factible original; por eso,
aceptar mediante capacidad agregada o caja bilateral tampoco prueba soporte ni
fricción.

El brazo $d$ actúa en sentidos opuestos en N1 y N2. Aumentarlo facilita el
momento de \eqref{eq:n2-pilot-demand}, mientras una formación más extendida
reduce la envolvente cinemática. N2 modifica la dispersión geométrica, pero
mantiene el número de miembros; por esa razón no valida
\eqref{eq:opposite-monotonicity}.

% --------------------------------------------------------------------------
% Página 7 · N3 formulación
% --------------------------------------------------------------------------
\leveltag{ESTIMACIÓN LOCAL DEL ESTADO Y DEL CONTACTO}
\section*{Pioneer 3-DX: SLAM, pose relativa y contacto}

El manual del fabricante documenta una base diferencial con motores CC
reversibles, encoders ópticos en cuadratura y sonar
\parencite{activmedia2003pioneer}. El manual no documenta LiDAR, IMU, cámara ni
sensor de fuerza para la unidad disponible; la arquitectura los incorpora como
instrumentación añadida que deberá calibrarse. Encoders y sonar pertenecen a la
base documentada; LiDAR, IMU, cámara, fiduciales y medida de contacto son una
cadena propuesta, todavía no integrada \parencite{olson2011apriltag}.


\paragraph{Transferencia instrumentada al Pioneer 3-DX.}
Una transferencia física exigiría combinar el RBPF con encoder, IMU y LiDAR de
\textcite{grisetti2007gridMapping}, un grafo relativo en $SE(2)$
\parencite{howard2005multirobotSlam,carlone2011linearGraphSlam}, fiduciales entre
base y carga \parencite{olson2011apriltag} y una confirmación de contacto
calibrada. Esa cadena no se ejecuta en N3: la campaña recibe estados sintéticos
y compara líder--seguidor con estructura virtual
\parencite{desai2001formation,ren2004virtual}. El primero retransmite el estado
por un árbol y el segundo actualiza copias vecinales por consenso
\parencite{kia2019dynamicConsensus,bullo2009distributed}; por tanto, N3 solo
mide calidad informativa y coste de red.

% --------------------------------------------------------------------------
% Página 8 · N3 experimento
% --------------------------------------------------------------------------
\leveltag{CALIDAD DE ESTIMACIÓN Y COSTE DE LA RED LOCAL}
\section*{Retardo, degradación, alcance y ancho de banda}

Cada condición de red se ejecuta con cuatro AMR, 80 actualizaciones y 30
semillas. El canal degradado combina pérdida $0{,}25$, una retransmisión y
$5$ kB/s, por lo que el diseño no aísla el efecto de la pérdida. Se registran tráfico,
latencia, edad y RMSE con ruido común de pose
$(0{,}020\ \mathrm m;\ 0{,}020\ \mathrm m;\ 0{,}010\ \mathrm{rad})$. El canal es
abstracto: no ejecuta SLAM, fusión IMU--encoder ni fiduciales; el estado
verdadero solo puntúa el error.

La Figura~\ref{fig:sp2-n3-comm} informa por separado el error de estimación, la
edad del estado y el coste de comunicación. Un intervalo que cruza cero impide
afirmar una diferencia.
\resultfigure[0.30\textheight]{figures/generated/fig-sp2-n3-quality-communication.pdf}{Bytes entregados, RMSE, diferencia pareada, edad y razón de entrega bajo cuatro condiciones de red; las barras del contraste indican IC del 95\%.}{fig:sp2-n3-comm}

\textbf{Resultado N3.} Transmitir más información no mejora siempre la
estimación: bajo el canal degradado, los mensajes adicionales envejecen en la
cola y la ventaja cambia de signo. En \SPTwoReadyNThreeRows{} ejecuciones,
líder--seguidor y estructura virtual
ofrecen en promedio \SPTwoReadyNThreeLeaderBytes{} y
\SPTwoReadyNThreeVirtualBytes{} bytes, de los que se entregan
\SPTwoReadyNThreeLeaderDeliveredBytes{} y
\SPTwoReadyNThreeVirtualDeliveredBytes{};
sus RMSE agregados son
$(\spdecimal{\SPTwoReadyNThreeLeaderPositionRMSE}\ \mathrm m,
\spdecimal{\SPTwoReadyNThreeLeaderYawRMSE}\ \mathrm{rad})$ y
$(\spdecimal{\SPTwoReadyNThreeVirtualPositionRMSE}\ \mathrm m,
\spdecimal{\SPTwoReadyNThreeVirtualYawRMSE}\ \mathrm{rad})$. Bajo el canal
degradado, la diferencia virtual menos líder es
$\spdecimal{\SPTwoReadyNThreeLossyPositionDifference}$ m, IC del 95\%
$[\spdecimal{\SPTwoReadyNThreeLossyPositionCILow},
\spdecimal{\SPTwoReadyNThreeLossyPositionCIHigh}]$, lo que refuta la mejora
uniforme. En ese canal, las razones de entrega son
$\spdecimal{\SPTwoReadyNThreeLossyLeaderDelivery}$ y
$\spdecimal{\SPTwoReadyNThreeLossyVirtualDelivery}$, casi iguales, pero la edad
media en línea pasa de $\spdecimal{\SPTwoReadyNThreeLossyLeaderAge}$ s a
$\spdecimal{\SPTwoReadyNThreeLossyVirtualAge}$ s. La estructura virtual ofrece
640 mensajes frente a 240 y su cola a 5 kB/s acumula estados obsoletos; esa
diferencia protocolaria explica el cambio de signo. Los datos cubren cuatro AMR
y un anillo; no caracterizan otros tamaños o grafos.

% --------------------------------------------------------------------------
% N4 · Controladores y guarda muestreada
% --------------------------------------------------------------------------
\leveltag{CONTROL NOMINAL Y GUARDA MUESTREADA DE ADMISIBILIDAD}
\section*{Seis leyes nominales bajo una comprobación común}

El estado $x=(q_L,\xi_L)$ se integra por Euler semiimplícito con
$\Delta t=0{,}20$ s. La dinámica de \eqref{eq:load-dynamics} convierte la
aceleración nominal $a_L^0$ en wrench. Las ganancias se fijaron manualmente
antes de abrir las semillas; no hubo conjunto de ajuste disjunto ni optimización
independiente por familia. Por ello, la comparación entre leyes es descriptiva.
La inferencia contrasta cada ley consigo misma, con y sin guarda.

Las ganancias congeladas de las leyes causales fueron
\[
\begin{aligned}
K_p^{\mathrm{PD}}&=(1{,}6;1{,}6;2{,}2), &
K_d^{\mathrm{PD}}&=(2{,}4;2{,}4;2{,}2),\\
K_i^{\mathrm{PID}}&=(0{,}12;0{,}12;0{,}08), &
K_s^{\mathrm{RA}}&=(1{,}9;1{,}9;2{,}7),\\
&&K_d^{\mathrm{RA}}&=(2{,}8;2{,}8;2{,}6).
\end{aligned}
\]
LQR usa $Q_L=\operatorname{diag}(8;8;10;3;3;3)$ y $R=I$. En LQI,
$Q_I=\operatorname{diag}(q_I)$ con
$q_I=(10;10;12;3;3;3{,}5;0{,}7;0{,}7;0{,}5)$. El MPC central, con horizonte 8,
usa $Q_M=\operatorname{diag}(8;8;10;2{,}2;2{,}2;2{,}6)$,
$R_M=\operatorname{diag}(0{,}35;0{,}35;0{,}4)$ y cotas
$(1{,}2;1{,}2;1{,}0)$. La Tabla~\ref{tab:sp2-methods} resume el papel y el
límite de cada familia.

Las seis leyes regulan la pose de la carga con geometría de contacto fija; no
controlan el error de formación de cada AMR. Defínase
\begin{equation}
\mathcal A_{\mathcal C}(x)=\left\{a:
\xi_L+\Delta t\,a\in\Xi_{\mathcal C}^{\mathrm{align}},\quad
W_d(x,a)\in\widehat{\mathcal W}_{\mathcal C}\right\}.
\label{eq:n4-admissible-acceleration}
\end{equation}
Las cinco leyes no predictivas usan $\widehat x_{L,\ell}$, la última estimación
del AMR de referencia; el MPC recibe $x_L$ global y actúa como techo central. El
consenso solo alinea copias de la pose. No se ejecutan DMPC, control distribuido
de formación ni reparto vecinal de wrench. PID y LQI congelan la integral cuando
interviene la guarda.

En S5--S6, un filtro inspirado en CBF exige
$n^\top a\geq-\alpha_1n^\top v-\alpha_0h$, donde
$h=\|p_L-p_o\|-r_{\mathcal C}-r_o-d_{\mathrm{safe}}$
\parencite{ames2017cbf}. La linealización omite la curvatura de $\ddot h$; por
tanto, solo certifica el candidato de la muestra actual. El gobernador recorre
$\mathcal A_\alpha=\{1,0{,}8,0{,}6,0{,}4,0{,}2\}$ y elige
\begin{equation}
\alpha^\star=\max\left\{\alpha\in\mathcal A_\alpha:
\alpha\widetilde a_L^0\in\mathcal A_{\mathcal C}(x),\quad
g_{\mathrm{CBF}}(x,\alpha\widetilde a_L^0)\geq0,\quad
g_{\mathrm{info}}=1\right\}.
\label{eq:n4-governor}
\end{equation}
Si el conjunto es vacío o la edad supera $\bar\tau=0{,}60$ s, se prueba
$a_L^{\mathrm{br}}=-1{,}8\xi_L$ con las mismas escalas y comprobaciones.
\textsc{Hold} solo se admite con velocidad pequeña. Si tampoco existe frenado
certificable, el modo es \textsc{uncontrolled}: el cero comandado no garantiza
detener una carga que ya se mueve y se requiere una parada externa.

\leveltag{SIETE ESCENARIOS Y COMPARACIÓN PAREADA}
\section*{Perturbaciones, fallo y criterio de éxito}

El diseño contiene siete escenarios y 30 semillas: \SPTwoReadyWorlds{} mundos
pareados y \SPTwoReadyDynamicRuns{} ejecuciones. S0 es nominal; S1 aumenta la
carga a 28 kg e introduce un impulso lateral; S2 heterogeneiza ruedas, tracción
y fricción; S3 impone pérdida $0{,}25$, retardo medio $0{,}16$ s y
$3{,}5$ kB/s; S4 retira un contacto a 3 s y lo restaura exógenamente a 4,5 s;
S5 añade obstáculo; S6 combina carga de 24 kg, pérdida permanente, red
degradada, obstáculo y perturbación. Se usan cuatro AMR, uno por esquina; el
tamaño de la coalición no es un factor del experimento.

En 12 s, el éxito requiere mantener durante $0{,}80$ s
\[
\|e_p\|<0{,}10\ \mathrm m,\quad |e_\theta|<0{,}10\ \mathrm{rad},\qquad
\|v_L\|<0{,}12\ \mathrm{m/s},\quad |\omega_L|<0{,}10\ \mathrm{rad/s}.
\]
Colisión, \emph{timeout} y
fallo numérico cuentan como fracaso. La huella mide
$r_{\mathcal C}=0{,}721$ m. La unidad inferencial es la semilla: cada remuestreo
conserva sus siete escenarios y las cuatro leyes con ablación. Los 210 pares
escenario--semilla no se tratan como observaciones independientes. Los IC usan
5.000 remuestreos de 30 bloques, Wilcoxon sobre la diferencia media por semilla
y corrección de Holm. La planta añade perturbación y ruido, pero fija contactos
y omite deformación y dinámica de rueda: es simulación reducida, no validación
física.

La Figura~\ref{fig:sp2-master-n4} concentra el patrón por escenario, el terminal
dominante, los contrastes pareados y dos medidas de admisibilidad realizada.
\resultfigure[0.59\textheight]{figures/generated/fig-sp2-master-n4.pdf}{Benchmark N4. A: éxito geométrico de las seis leyes con guarda. B: terminal dominante; AGE, información obsoleta; KIN, límite de ruedas; SUP, soporte inviable; COL, colisión. C: efecto guarda menos crudo e IC del 95\% por 30 bloques. D: fracción y duración de violaciones realizadas para las cuatro leyes pareadas.}{fig:sp2-master-n4}

\textbf{Patrón por escenario.} Las seis leyes alcanzan el terminal en S0, S2,
S4 y S5; S1 solo registra fallos con PID. S3 separa las leyes locales del MPC
central: el éxito agregado baja a
$\spdecimal{\SPTwoReadySThreeSuccess}$. En S5 no se observan colisiones, pero
esa muestra no demuestra llegada general ni invariancia continua. S6 produce el
resultado negativo principal: éxito
$\spdecimal{\SPTwoReadySSixSuccess}$ y colisión
$\spdecimal{\SPTwoReadySSixCollision}$. Tras la pérdida permanente, tres
contactos no sostienen la carga. La barrera no puede reponer el apoyo ausente ni
prometer frenado cuando el conjunto admisible queda vacío.

\textbf{Efecto pareado.} La guarda cambia el éxito geométrico en
$+\spdecimal{\SPTwoReadyPairedSuccessDifference}$, IC del 95\%
[$+\spdecimal{\SPTwoReadyPairedSuccessCILow}$;
$+\spdecimal{\SPTwoReadyPairedSuccessCIHigh}$]. Para el éxito físicamente
admisible según el sustituto 2.5D, el efecto es
$+\spdecimal{\SPTwoReadyPairedAdmissibleSuccessDifference}$, IC
[$+\spdecimal{\SPTwoReadyPairedAdmissibleSuccessCILow}$;
$+\spdecimal{\SPTwoReadyPairedAdmissibleSuccessCIHigh}$]. El efecto terminal es
positivo para PD, LQI y amortiguamiento y nulo para MPC.

H6 queda \textbf{refutada}: la diferencia agregada del riesgo de colisión es
$+\spdecimal{\SPTwoReadyPairedCollisionDifference}$, IC
[$+\spdecimal{\SPTwoReadyPairedCollisionCILow}$;
$+\spdecimal{\SPTwoReadyPairedCollisionCIHigh}$]. El desglose por ley es
diagnóstico \emph{post hoc}: PD, LQI y amortiguamiento aumentan el riesgo en
$0{,}124$, $0{,}114$ y $0{,}114$; MPC lo reduce en $0{,}143$. En estos mundos,
mejorar la admisibilidad instantánea no equivale a mejorar la seguridad de la
trayectoria.

La fracción temporal media fuera del modelo baja de
$\spdecimal{\SPTwoReadyPairedRawRealizedFraction}$ a
$\spdecimal{\SPTwoReadyPairedGovernedRealizedFraction}$; el episodio continuo
más largo, de $\spdecimal{\SPTwoReadyPairedRawMaxViolationDuration}$ s a
$\spdecimal{\SPTwoReadyPairedGovernedMaxViolationDuration}$ s. La sensibilidad
varía $\Delta\alpha$, no $\Delta t$: no estudia discretización temporal ni
establece invariancia continua. El resultado exporta una interfaz operativa en
el modelo plano, no un certificado de estabilidad híbrida ni de contacto real.

% --------------------------------------------------------------------------
% Demostrador Cargo integrado · evidencia existente y auditada
% --------------------------------------------------------------------------
\leveltag{MISIÓN CARGO INTEGRADA HASTA LA POSE DE ENTREGA}
\section*{Aproximación, transporte y sustitución}

La campaña enlaza reclutamiento, llegada posicional de uniciclos, transporte
planar, fallo, viaje de un sustituto y reanudación. La carga permanece inmóvil
mientras el certificado agregado está abierto. La Figura~\ref{fig:sp2-master-cargo}
combina la traza espacial disponible, las secuencias de fase observadas y la
ablación por tratamiento. No reconstruye poses de AMR ni tiempos de transición:
esos datos no se almacenaron. La misión termina al alcanzar la pose de entrega;
la liberación física no se ejecuta.

\resultfigure[0.57\textheight]{figures/generated/fig-sp2-master-cargo.pdf}{Misión Cargo integrada. A: pose de la carga en un mundo nominal representativo; la campaña no archivó trayectorias individuales de AMR. B: secuencias categóricas registradas, sin interpretarlas como eje temporal; \textsc{Recuperar} comprende llegada de reserva y revalidación lógica, no contacto físico. C: tasa de misión por tratamiento y régimen.}{fig:sp2-master-cargo}

La arquitectura es híbrida. Los vecinos intercambian identificadores, pero el
líder, los atributos, la selección espacial, la pose de carga y el reparto de
wrench consultan estado o registro global. El \emph{docking} solo verifica
llegada a $0{,}12$ m y velocidad inferior a $0{,}08$ m/s durante cinco muestras;
no valida orientación, contacto ni colisiones de aproximación. Después, la
geometría de los AMR se impone rígidamente respecto de la carga.

El diseño contiene cuatro regímenes, $N\in\{4,8,12\}$, 30 semillas,
\CargoEtwoEWorlds{} mundos y \CargoEtwoERuns{} ejecuciones. Solo el régimen de
red degradada usa diez eventos, pérdida $0{,}25$ y retardo máximo de dos; los
otros usan siete eventos, pérdida $0{,}02$ y retardo nulo. El método vecinal
completo alcanza la entrega en 359/360 mundos, las 360 llegadas posicionales y
89/90 sustituciones; la referencia central completa 360/360.

Los contrastes confirmatorios matizan esa tasa. E2E-H2 respalda el bloque
mecánico sin aislar sus componentes: efecto $0{,}996$, IC del 95\%
$[0{,}989;1{,}000]$, $p_{\mathrm{Holm}}=5{,}79\times10^{-21}$. E2E-H3 respalda
reparar frente a detenerse: efecto $0{,}989$, IC
$[0{,}967;1{,}000]$, $p_{\mathrm{Holm}}=6{,}46\times10^{-27}$. E2E-H1 no muestra
ventaja temporal ($p_{\mathrm{Holm}}=0{,}068$) y E2E-H4 no prueba equivalencia
con información perfecta. La campaña acredita una misión simulada hasta la pose
de entrega, no percepción, middleware, contacto 3D, liberación ni hardware.


% --------------------------------------------------------------------------
% Página 9 · Extensión caging, Figura 5 conservada
% --------------------------------------------------------------------------
\leveltag{ESTABILIDAD NOMINAL DEL SEGUIMIENTO RÍGIDO}
\section*{Qué garantiza el lazo y qué lo condiciona}

Las dos secciones anteriores acotan cuándo una orden es cinemáticamente
admisible y cuándo el wrench pedido es realizable. Falta decir qué ocurre con el
error de pose cuando ambas condiciones se cumplen. Sea
$e=(x_L-x_d,\;y_L-y_d,\;\operatorname{wrap}(\theta_L-\theta_d))$ el error frente
a una referencia dos veces derivable y
\begin{equation}
W_d^0=M_La_d^0+D_L\xi_L-K_D\dot e-K_Pe,
\qquad K_P\succ0,\ K_D\succ0 .
\label{eq:computed-wrench}
\end{equation}

\begin{sp2prop}[Convergencia del error de pose]
Si la referencia se mantiene en la envolvente cinemática de la coalición, el
reparto de fuerzas realiza exactamente el wrench de \eqref{eq:computed-wrench},
los parámetros $M_L$ y $D_L$ son los del modelo y no hay saturación ni
perturbación no modelada, entonces $e\to0$ y $\dot e\to0$.
\end{sp2prop}

Sustituyendo \eqref{eq:computed-wrench} en \eqref{eq:sp2-payload-el} queda
$M_L\ddot e+K_D\dot e+K_Pe=0$. Con
$V=\tfrac12\dot e^\top M_L\dot e+\tfrac12e^\top K_Pe$ se obtiene
$\dot V=-\dot e^\top K_D\dot e\leq0$, y el mayor conjunto invariante contenido
en $\{\dot V=0\}$ es el origen. La envoltura de ángulo hace que el resultado sea
local: excluye el conjunto de errores de rumbo de magnitud $\pi$.

Las cuatro hipótesis no son formalidades, y dos de ellas las mide este capítulo.
La primera falla en \SPTwoReadyNOneFalse{} de las \SPTwoReadyNOneRows{} órdenes
del barrido, que el comparador por pivote aceptaría y la comprobación conjunta
rechaza. La segunda falla en \SPTwoReadyNTwoFalse{} de \SPTwoReadyNTwoRows{}
casos, en los que el reparto no realiza el wrench dentro de los conos de
fricción. El
lazo no es incondicionalmente estable: es estable dentro de una envolvente cuyo
tamaño se ha medido, y los experimentos N1 y N2 son precisamente la medida de
cuántas veces se sale de ella. Los parámetros dinámicos del Pioneer no están
identificados, así que la tercera hipótesis queda declarada y no verificada.

\leveltag{CIERRE GEOMÉTRICO BAJO CONTACTO UNILATERAL}
\section*{Verificador traslacional en cuadrícula}

El confinamiento no prehensil sigue una relación de contacto distinta de Cargo
\parencite{pereira2004caging,fink2008caging}, formalizada en
\eqref{eq:sp2-caging-complementarity} y \eqref{eq:sp2-caging-el}. El verificador
de esta sección trabaja sobre la separación $g_i^{E}$ y no sobre las reacciones:
comprueba geometría de escape, no equilibrio de fuerzas. Sea $\mathcal F_{\delta_{xy}}$ el
conjunto de celdas libres y $\operatorname{Comp}_{\delta_{xy}}(q_0)$ la
componente alcanzable desde $q_0$ mediante cuatro vecinos. El cierre se
certifica mediante
\begin{equation}
\operatorname{conf}_{\delta_{xy}}(q_0)=1
\quad\Longleftrightarrow\quad
\operatorname{Comp}_{\delta_{xy}}(q_0)
\cap\partial\mathcal F_{\delta_{xy}}=\varnothing .
\label{eq:grid-cage}
\end{equation}


Con $21\times21$ celdas y $\delta_{xy}=0{,}05$ m, el verificador cierra el anillo
y detecta la abertura, pero omite orientación, fricción, modos y fuerzas
\parencite{pereira2004caging,fink2008caging}. Como referencias dinámicas,
\textcite{lee2025switching} emplean conmutación central MILP--QP y
\textcite{rosenfelder2024force}, control de fuerza; el DMPC vecinal no se
implementa aquí.

La Tabla~\ref{tab:sp2-evidence} separa el tamaño de cada campaña de la
afirmación que permite sostener. Esta lectura evita tratar una simulación
cinemática, un contraste pareado y un verificador geométrico como evidencias
intercambiables.
\begin{table}[H]
\centering
\fontsize{8.4}{8.95}\selectfont
\setlength{\tabcolsep}{3.0pt}
\renewcommand{\arraystretch}{0.94}
\caption{Resultados de SP2, unidad experimental y límite de validez.}
\label{tab:sp2-evidence}
\begin{tabularx}{\textwidth}{p{0.90cm} p{2.45cm} p{5.20cm} X}
\toprule
Nivel & Unidad y tamaño & Resultado observado & Alcance acreditado \\
\midrule
N1 & 36.600 órdenes & 4.536 falsos positivos del criterio simple (12{,}4\,\%) & Cinemática de pivote y ruedas; sin par ni contacto \\
N2 & 720 demandas & 668/720 factibles con soporte y fricción & Sustituto 2.5D con contactos fijos \\
N3 & 240 ejecuciones; 30 semillas & En red degradada, $\Delta\mathrm{RMSE}=0{,}132$ m, IC [0{,}131; 0{,}134] & Cuatro AMR en anillo; canal abstracto, no SLAM ejecutado \\
N4 & 210 mundos; 30 bloques & Guarda: $\Delta P_{\rm éxito}=+0{,}067$; $\Delta P_{\rm col}=+0{,}052$ & Planta reducida; H6 de colisión refutada \\
Cargo & 360 mundos; 2.160 ejecuciones & Vecinal completo 359/360; sustitución 89/90 & Hasta pose de entrega; arquitectura híbrida y acoplamiento posicional \\
Caging & verificador $21\!\times\!21$ & Cierre o ruta de escape por BFS de cuatro vecinos & Certificado geométrico traslacional; sin orientación, fuerza ni dinámica \\
\bottomrule
\end{tabularx}

\end{table}

% --------------------------------------------------------------------------
% Página 12 · Síntesis y transición a SP3
% --------------------------------------------------------------------------
\leveltag{CONCLUSIÓN DE SP2 Y ENVOLVENTE PARA LA PLANIFICACIÓN}
\section*{Resultado exportado e interfaz con SP3}

La cota de \eqref{eq:pure-rotation-bound} y la
Figura~\ref{fig:sp2-master-n1-n2} impiden que SP3 trate la coalición como un
punto. Su envolvente computada es
\begin{equation}
\widehat{\mathcal U}_{\mathcal C}=
\{u:\xi_L(u)\in\Xi_{\mathcal C}^{\mathrm{align}},\ W_d(u)\in
\widehat{\mathcal W}_{\mathcal C},\ g_{\mathrm{CBF}}\geq0,\ g_{\mathrm{info}}=1\}.
\label{eq:virtual-vehicle}
\end{equation}
El contrato conserva rejilla, margen, huella, $\mathcal A_\alpha$, terminal y
modos. N1 fija la equivalencia
cinemática; N2 encuentra \SPTwoReadyNTwoFalse{} falsos positivos; N3 descarta
una mejora informativa uniforme; N4 mejora terminal y admisibilidad, pero
refuta H6; Cargo llega en 359/360 misiones y sustituye en 89/90. No se validan
acoplamiento físico, liberación, estabilidad, SLAM, fuerza ni 3D; caging queda
geométrico. La Tabla~\ref{tab:sp2-evidence} conserva esa frontera al exportar el
resultado a SP3.

\end{document}
